\scanpage{4}

\mspage{\uline{2}}

\deleted{\diagram{\begin{tikzcd}[ampersand replacement=\&,column sep=small,row sep=small]
\tilde x \arrow[r,"{\deleted{\unsure{\varphi}}}"] \arrow[rr,bend left,"\varphi"] \arrow[d,"\sim"'] \& \tilde y \arrow[r,tail,"i","\sim"'] \arrow[d,"\sim"] \& \tilde y' \arrow[rr,bend left,"\psi"] \& \tilde z \arrow[r,tail,"j","\sim"'] \arrow[d,"\sim"'] \& \tilde z' \arrow[dl] \\
x \arrow[r,"f",no head] \& y \arrow[rr,"g"] \& {} \& z \& {}
\end{tikzcd}}}

\deleted{\diagram{\begin{tikzcd}[ampersand replacement=\&,column sep=small,row sep=small]
\tilde x \arrow[dr] \& \tilde y \arrow[d,"i"] \arrow[dr,"\psi"] \& \tilde z \arrow[d,"j"] \\
{} \& \tilde y' \& \tilde z'
\end{tikzcd}}}

\diagram{\begin{tikzcd}[ampersand replacement=\&,column sep=small,row sep=small]
\tilde x \arrow[r,"{\deleted{\unsure{\tilde f|\tilde\varphi}}}"] \arrow[dr,"{\tilde f}"'] \arrow[ddd,dashed,"\sim"'] \& \tilde y \arrow[d,tail,"{i_f}"] \arrow[dr,"{\tilde g}"] \& \tilde z \arrow[d,tail,"{i_g}"] \\
{} \& \tilde y' \arrow[dd,dashed,"\sim"'] \arrow[dr,no head,"{\deleted{\unsure{u}}\tilde g'}"] \& \tilde z' \arrow[d,tail] \\
{} \& {} \& \tilde z'' \arrow[d,dashed,no head] \\
x \arrow[r,"f"] \& y \arrow[r,"g"] \& z
\end{tikzcd}}

\deleted{\unsure{On|Ce}} Montrons que ce foncteur transforme $f\in W$ en iso en effet, si $f\in W$, donc $f\alpha_x\in W$, on voit aussitôt qu'on aura $\widetilde f\in W$, donc $[\widetilde f]$ est un iso.

Considérons les deux composés

\diagram{\begin{tikzcd}[ampersand replacement=\&,column sep=small,row sep=small]
W_{\rm Cof}^{-1}\,\mathrm{Cof} \arrow[r,shift left,"\Psi"] \& W^{-1}C \arrow[l,shift left,"\Phi"]
\end{tikzcd}}

\deleted{pour voir} \inserted{prouvons} qu'\deleted{ils sont} sont \deleted{\unsure{égaux}} iso aux foncteurs identiques.

Pour $\Phi\Psi$, c'est $=$ l'identité, si on prend soin de choisir $\widetilde x\xrightarrow{\alpha_x}x$ l'identité, quand $x\in\mathrm{Cof}$. Pour $\Psi\Phi$, il \deleted{suffit de} \inserted{\unsure{veut|vaut|faut}} voir que le foncteur \deleted{$\Psi\Phi\overset{\rm déf}{=}$} $R\overset{\rm déf}{=}\Psi\Phi\mid C$

\[
 x\longmapsto\widetilde x\qquad C\longrightarrow W^{-1}C
\]

est iso \deleted{à l'identité} \inserted{au foncteur can}, on a

\diagram{\begin{tikzcd}[ampersand replacement=\&,column sep=small,row sep=small]
R(x) \arrow[d,equal] \arrow[r,"{[\alpha_x]}"] \& x \\
\widetilde x \& {}
\end{tikzcd}}

dans $C$ donc dans $W^{-1}C$ (soit $[\alpha_x]$ qui est un iso, il faut vérifier que c'est fonctoriel

\diagram{\begin{tikzcd}[ampersand replacement=\&,column sep=small,row sep=small]
\widetilde x \arrow[d,"{\alpha_x}"'] \arrow[rr,bend left,"{\widetilde f}"] \& \widetilde y \arrow[r,tail,"{i_f}"'] \arrow[d,"{\alpha_y}"] \& \widetilde y' \arrow[dl,"{\beta_f}"] \\
x \arrow[r,"f"] \& y \& {}
\end{tikzcd}}

i.e qu'on a
\[
 [\alpha_y]\deleted{\unsure{{}\circ}}\bigl([i_f]^{-1}\deleted{\unsure{{}[\,]}}[\tilde f]\bigr)=[f][\alpha_x]
\]
Mais $\alpha_y=\beta_f i_f$ d'où le premier mb est $[\beta_f\tilde f]=[f\alpha_x]=[f][\alpha_x]$ cqfd.
% archive sheet number: 3
